\documentclass[10pt,letterpaper]{article}

\usepackage[letterpaper,top=0.34in,bottom=0.32in,left=0.40in,right=0.40in]{geometry}
\usepackage{fontspec}
\setmainfont{DejaVu Sans}
\usepackage{xcolor}
\usepackage{enumitem}
\usepackage{hyperref}
\usepackage{titlesec}
\usepackage{microtype}

\definecolor{nameblue}{HTML}{214F78}
\definecolor{linkblue}{HTML}{0000CC}

\hypersetup{
  colorlinks=true,
  urlcolor=linkblue,
  linkcolor=linkblue,
  pdfauthor={Harsh Desai},
  pdftitle={Harsh Desai Resume - UW IDP Mechanical Instrumentation}
}

\pagestyle{empty}
\raggedright
\setlength{\parindent}{0pt}
\setlength{\tabcolsep}{0pt}
\setlength{\emergencystretch}{1em}
\hyphenpenalty=10000
\exhyphenpenalty=10000
\fontsize{8.40pt}{9.05pt}\selectfont

\titleformat{\section}
  {\fontsize{9.50pt}{9.95pt}\selectfont\bfseries\uppercase}
  {}{0pt}{}
  [\vspace{-1.8pt}\titlerule]
\titlespacing*{\section}{0pt}{3.9pt}{1.2pt}

\setlist[itemize]{
  leftmargin=13pt,
  label=\raisebox{0.17ex}{\tiny$\bullet$},
  itemsep=0.15pt,
  parsep=0pt,
  topsep=0.20pt,
  partopsep=0pt
}

\newcommand{\resumeSubheading}[4]{%
  \par\noindent\begin{tabular*}{\textwidth}{@{\extracolsep{\fill}}l r}
    \textbf{#1} & \textcolor{nameblue}{\textbf{#2}} \\
    #3 & #4
  \end{tabular*}\par\vspace{-1.0pt}%
}

\newcommand{\resumeProjectHeading}[2]{%
  \par\noindent\begin{tabular*}{\textwidth}{@{\extracolsep{\fill}}l r}
    \textbf{#1} & #2
  \end{tabular*}\par\vspace{-1.0pt}%
}

\newcommand{\resumeItemsStart}{\begin{itemize}}
\newcommand{\resumeItemsEnd}{\end{itemize}\vspace{-0.5pt}}
\newcommand{\resumeItem}[1]{\item #1}

\begin{document}

\begin{tabular*}{\textwidth}{@{\extracolsep{\fill}}l r}
  {\fontsize{18pt}{18pt}\selectfont\bfseries\textcolor{nameblue}{Harsh Desai}} &
  \fontsize{8.0pt}{8.6pt}\selectfont
  Ann Arbor, MI $\mid$
  \href{mailto:harshdes@umich.edu}{harshdes@umich.edu} $\mid$
  +1-734-548-1080 $\mid$
  \href{https://www.linkedin.com/in/harshddes/}{LinkedIn} $\mid$
  \href{https://harshddes.github.io/}{Portfolio}
\end{tabular*}
\vspace{-3.0pt}

Mechanical and Space Systems Engineer experienced in CAD/CAE, electromechanical mechanisms, scientific instrumentation, vacuum/HV test stands, subsystem integration, and design validation.

\section{Education}
\resumeSubheading
  {University of Michigan}{Ann Arbor, MI}
  {Master of Engineering, Space Systems Engineering $\mid$ GPA: 3.79/4.0}{Aug 2024-Dec 2025}
\resumeSubheading
  {Vellore Institute of Technology}{Vellore, India}
  {Bachelor of Technology, Mechanical Engineering $\mid$ GPA: 7.78/10}{Jul 2019-Jul 2023}

\section{Technical Skills}
\textbf{Mechanical System Design \& CAE:} SolidWorks, Fusion 360, Autodesk Inventor, AutoCAD, ANSYS Mechanical/Fluent, SpaceClaim; conceptual design, 3D part/assembly modeling, structural and thermal-fluid analysis\\[-0.2pt]
\textbf{Prototype \& Test:} mechanism integration, component selection and procurement, assembly, troubleshooting, test procedures, design verification, FMEA, vacuum-chamber operation, high-voltage test stands\\[-0.2pt]
\textbf{Instrumentation \& Automation:} plasma diagnostics, sensors and detector systems, electrostatic ion optics, Python, PyVISA/PyMeasure, LabVIEW (CLAD preparation), MATLAB, Keithley DAQs/SMUs, TDK Lambda power supplies\\[-0.2pt]
\textbf{Documentation \& Analysis:} Microsoft Excel, Word, PowerPoint; test procedures, FMEA, requirements and verification documentation

\section{Engineering Experience}
\resumeSubheading
  {Climate and Space Sciences and Engineering, University of Michigan}{Ann Arbor, MI}
  {Research Assistant, LVACCS Experimental Hardware \& Test Automation}{Feb 2026-Present}
\resumeItemsStart
  \resumeItem{Integrated and validated a 1.3~kV plasma-source test stand spanning vacuum hardware, a custom discharge interface, power supplies, Keithley acquisition, and synchronized Python/PyVISA control; reduced manual steps from 15 to 5}
  \resumeItem{Developed a discharge-box FMEA, ignition and acquisition procedures, and post-run checks; achieved 98.6\% synchronized data capture and improved repeatability across hardware test states}
  \resumeItem{Operated Langmuir and emissive probes plus an RPA assembly during vacuum-chamber testing, performing LP/EP measurements and synchronized RPA retarding-voltage/collector-current acquisition while resolving mechanical/electrical interface issues}
\resumeItemsEnd

\resumeSubheading
  {Space Physics Research Lab, University of Michigan}{Ann Arbor, MI}
  {Summer Intern, TestBedz Spacecraft Qualification Platform}{May 2025-Jul 2025}
\resumeItemsStart
  \resumeItem{Built a requirements-to-test workflow for thermal-vacuum, vibration, and EMI qualification, matching hardware to facility limits across six test types}
  \resumeItem{Supported subsystem design reviews and technical presentations by translating operating envelopes and acceptance criteria into compatibility checks and traceable test plans for repeatable handoff to test facilities}
\resumeItemsEnd

\resumeSubheading
  {Solar and Heliospheric Research Group, University of Michigan}{Ann Arbor, MI}
  {Graduate Research Assistant, Solid-State Detector Readout}{Jan 2025-Dec 2025}
\resumeItemsStart
  \resumeItem{Derived ADC requirements from detector energy resolution, noise, signal amplitude, and timing; evaluated a 14-bit, 125~MSPS ADC integration with a Zynq-7000 development platform}
  \resumeItem{Verified MATLAB HDL Coder models against Vivado co-simulation and diagnosed 125~MHz post-synthesis timing constraints, documenting the design limits and next integration steps}
\resumeItemsEnd

\section{Selected Mechanical \& Instrumentation Projects}
\resumeProjectHeading{CanSat 2022 (NASA/AAS) $\mid$ 7th worldwide of 42 teams}{Aug 2021-Jul 2022}
\resumeItemsStart
  \resumeItem{Designed and built a compact 10~m tether-deployment and camera-stabilization assembly using a DC motor, worm gear, bearings, an ABS printed structure, and a dual-servo gimbal for controlled descent and fixed 45-degree imaging}
  \resumeItem{Iterated the deployment architecture from spur to worm gearing for one-way load transfer, then worked across the mechanical and electronics team to integrate the mechanism with payload electronics and verify release and imaging functions through subsystem tests and design reviews}
\resumeItemsEnd

\resumeProjectHeading{Spaceport America Cup / IREC $\mid$ 23rd globally, 5th Asia-Pacific}{Apr 2020-Jul 2021}
\resumeItemsStart
  \resumeItem{Developed an in-house trajectory model for a Mach~0.9, 10,000~ft rocket, generating max-Q pressure and drag inputs in ANSYS Fluent and cross-checking the flight prediction in OpenRocket}
\resumeItemsEnd

\resumeProjectHeading{Particle Detector \& Electrostatic Analyzer Calibration}{Jul 2025-Dec 2025}
\resumeItemsStart
  \resumeItem{Calibrated a channel electron multiplier (CEM) signal chain by configuring the ion source, CEM, charge-sensitive preamplifier, pulse shaper, and MCA; charge-calibrated MCA channels with an RC injector and analyzed pulse-height spectra to extract centroid, FWHM, count rate, and effective gain versus detector bias}
  \resumeItem{Modeled a quadrispherical electrostatic analyzer in SIMION across an 84-point energy-angle grid, quantifying pass-energy shift, energy resolution, angular acceptance, and fringe-field sensitivity}
\resumeItemsEnd

\end{document}
